%----------------------------------------------------------------------------------------
%    PACKAGES AND THEMES
%----------------------------------------------------------------------------------------

\documentclass[aspectratio=169,xcolor=dvipsnames]{beamer}
\usetheme{DEXPI}
\logo{}
\usepackage[utf8]{inputenc}
\usepackage[T1]{fontenc}
\usepackage[portuguese]{babel}
\usepackage{hyperref}
\usepackage{graphicx}
\usepackage{booktabs}
\usepackage{xcolor}

% Configuração para áreas reservadas para imagens
\newcommand{\imageplaceholder}[1]{
    \begin{center}
        \fcolorbox{gray}{lightgray}{\begin{minipage}[c][0.35\textheight]{0.8\textwidth}
            \centering \textit{[Espaço para Imagem: #1]}
        \end{minipage}}
    \end{center}
}

%----------------------------------------------------------------------------------------
%    TITLE PAGE
%----------------------------------------------------------------------------------------

\title{NOX: Rumo a um Sistema Operacional para Redes}
\subtitle{O Paradigma SDN em Larga Escala}

\author{Apresentação de Seminário}
\institute{Programa de Pós-Graduação}
\date{\today}

\begin{document}

\begin{frame}[plain]
    \titlepage
\end{frame}

%------------------------------------------------
\section{O Desafio}
%------------------------------------------------

\begin{frame}{O Caos do Gerenciamento Tradicional}
    \textbf{Hook:} \textit{"Por que gerenciar redes continua tão difícil?"}
    
    \vspace{0.5cm}
    
    \begin{itemize}
        \item \textbf{Gestão Fragmentada:} Configuração manual de componentes isolados.
        \item \textbf{Políticas Engessadas:} Dependentes da topologia e do baixo nível (IPs, MACs).
        \item \textbf{Analogia:} Redes operam hoje como os computadores operavam antes da invenção do Sistema Operacional.
    \end{itemize}
    
    \vspace{0.3cm}
    \imageplaceholder{Esquema de uma rede confusa / emaranhado de protocolos}
\end{frame}

%------------------------------------------------
\section{A Proposta}
%------------------------------------------------

\begin{frame}{NOX: O Sistema Operacional de Redes}
    \textbf{Hook:} \textit{"E se a rede inteira pudesse ser controlada como uma única máquina?"}
    
    \vspace{0.5cm}
    
    \begin{itemize}
        \item \textbf{Visão Centralizada:} Uma única interface de programação para toda a infraestrutura.
        \item \textbf{Abstração de Nomes:} Operação sobre Identidades (Usuários) ao invés de Endereços.
        \item \textbf{OpenFlow:} O protocolo que permite o controle do roteamento no plano de dados.
    \end{itemize}
    
    \vspace{0.3cm}
    \imageplaceholder{Diagrama Controle vs. Plano de Dados (Controller central)}
\end{frame}

\begin{frame}{Arquitetura e Comportamento}
    \textbf{Hook:} \textit{"Toda a complexidade resolvida na análise do primeiro pacote."}
    
    \vspace{0.5cm}
    
    \begin{itemize}
        \item \textbf{Network View:} Um banco de dados que mantêm a topologia global atualizada.
        \item \textbf{Reação a Eventos:} Switches delegam os fluxos desconhecidos ao Controlador.
        \item \textbf{Programação Lógica:} As aplicações (ex: Firewall, Roteamento) rodam em cima do NOX, e não nos switches.
    \end{itemize}

    \vspace{0.3cm}
    \imageplaceholder{Ciclo de vida do processamento de pacotes (Switch -> Controlador -> Fluxo)}
\end{frame}

%------------------------------------------------
\section{Resultados Obtidos}
%------------------------------------------------

\begin{frame}{A Prova de Conceito: Resultados e Escalabilidade}
    \textbf{Hook:} \textit{"É possível ter controle centralizado sem sacrificar desempenho?"}
    
    \vspace{0.5cm}
    
    \begin{itemize}
        \item \textbf{Altíssimo Throughput:} +100.000 inícios de fluxo processados por segundo (Hardware Comum).
        \item \textbf{Drástica Redução de Código:} O projeto \textit{Ethane} foi reduzido de 45.000 linhas de C++ para poucas milhares de linhas de Python.
        \item \textbf{Homologado na Prática:} Único provedor de rede ativo em 30 hosts por +6 meses.
    \end{itemize}

    \vspace{0.3cm}
    \imageplaceholder{Gráfico comparativo de desempenho / linhas de código}
\end{frame}

%------------------------------------------------
\section{Conclusões}
%------------------------------------------------

\begin{frame}{Conclusões do Estudo}
    \textbf{Hook:} \textit{"O legado que o NOX deixou para o futuro das Redes (SDN)."}
    
    \vspace{0.5cm}
    
    \begin{itemize}
        \item \textbf{Fim do Paradigma Distribuído:} Escala massiva não requer algoritmos totalmente distribuídos.
        \item \textbf{Aceleração da Inovação:} Facilita desenvolvimento de redes sem se preocupar com limitações de hardware ou fabricantes.
    \end{itemize}

    \vspace{0.3cm}
    \imageplaceholder{Linha do tempo / transição de redes legadas para SDN}
\end{frame}

%------------------------------------------------
\section{Comentários Pessoais}
%------------------------------------------------

\begin{frame}{Visão Pessoal: Lições e Dificuldades}
    \textbf{Hook:} \textit{"Uma nova forma de pensar a conectividade."}
    
    \vspace{0.5cm}
    
    \begin{columns}[T]
        \begin{column}{0.48\textwidth}
            \textbf{Principais Lições}
            \begin{itemize}
                \item Desacoplar Inteligência (Controle) do Hardware (Dados) é revolucionário.
                \item Gerência tratada como Engenharia de Software pura.
            \end{itemize}
        \end{column}
        \begin{column}{0.48\textwidth}
            \textbf{Desafios e Dificuldades}
            \begin{itemize}
                \item Desconstruir o modelo mental dos protocolos convencionais (BGP/OSPF).
                \item Compreender a gerência do "Ponto único de falha" sob a lógica centralizada.
            \end{itemize}
        \end{column}
    \end{columns}

    \vspace{0.3cm}
    \imageplaceholder{Ícone ilustrando um desafio / barreira ultrapassada}
\end{frame}

%------------------------------------------------

\begin{frame}
    \Huge{\centerline{\textbf{Obrigado}}}
\end{frame}

\end{document}
